% from ega4-8-fr.tex
\oldpage[IV-3]{7}
«~Géométrie algébrique absolue~»). Cet exemple nous montre en particulier que
dans la plupart des questions «~relatives~» sur un préschéma de base $Y$, on
peut se ramener au cas où $Y$ est \emph{noethérien}.

On notera que dans cet exemple, au contraire des précédents, les morphismes
$S_\mu\to S_\lambda$ n'ont en général aucune propriété particulière de
régularité.
\end{enumerate}

Par la suite, quand nous appliquerons les résultats qui vont suivre à l'une
quelconque des trois situations particulières qui viennent d'être décrites,
nous nous dispenserons de reprendre en détail le procédé qui permet de faire
ces applications, nous contentant de renvoyer à ce qui précède.
\end{env}

\begin{env}[8.1.3]
\label{IV.8.1.3-fr}
Dans l'exemple \emph{a)} de (8.1.2), nous avons vu que si $Y$ est un
préschéma intègre de point générique $y$, $f:X\to Y$ un morphisme de
présentation finie, alors, si la fibre générique $f^{-1}(y)$ est propre sur
$\boldsymbol k(y)$, il y a un voisinage ouvert $U$ de $y$ tel que $f^{-1}(U)$
soit propre sur $U$~; il en résulte \emph{a fortiori} que pour tout $s\in U$,
$f^{-1}(s)$ est propre sur $\boldsymbol k(s)$. Il arrive qu'on ait besoin
d'une réciproque, assurant que si $f^{-1}(s)$ est propre sur
$\boldsymbol k(s)$ pour «~suffisamment~» de points $s\in U$, alors
$f^{-1}(y)$ est propre sur $\boldsymbol k(y)$. Par exemple, supposons que $X$
et $Y$ soient des préschémas algébriques sur un corps $k$ algébriquement clos
(on peut prendre pour $k$ le corps $\mathbf C$ des nombres complexes, pour
fixer les idées)~; on a parfois besoin de savoir que si, pour tout $s\in Y$,
\emph{rationnel sur $k$}, la fibre $f^{-1}(s)$ est propre sur
$\boldsymbol k(s)$, alors $f^{-1}(y)$ est propre sur $\boldsymbol k(y)$, et
par suite $f^{-1}(U)$ est propre sur $U$ pour un voisinage $U$ de
$y$\footnote{On notera qu'un tel énoncé est en fin de compte purement
géométrique, en ce sens qu'il ne fait appel qu'à des points rationnels sur
$k$, et non à des points génériques~; par exemple, lorsque $k=\mathbf C$, cet
énoncé a une signification topologique évidente pour l'analyste, en
interprétant «~propre~» au sens topologique du terme, pour les espaces
sous-jacents aux espaces analytiques formés des points de $X$ et $Y$
rationnels sur $\mathbf C$.}. Or cet énoncé résultera aisément du suivant~:
l'ensemble $E$ des points $s\in Y$ tels que $f^{-1}(s)$ soit propre sur
$\boldsymbol k(s)$ est \emph{constructible} (et par conséquent identique à
$Y$ tout entier s'il contient les points fermés de $Y$, grâce au théorème des
zéros de Hilbert (10.4.8))~; cela équivaut encore à dire que si $f^{-1}(y)$
\emph{n'est pas propre} sur $\boldsymbol k(y)$, alors il existe un voisinage
ouvert $U$ de $y$ tel que $f^{-1}(s)$ ne soit pas propre sur
$\boldsymbol k(s)$ pour tout $s\in U$ (cf. (9.6.1, (iv))). Cet exemple
illustre l'intérêt qu'il y a à développer systématiquement des \emph{critères
de constructibilité} pour les notions les plus importantes~: c'est ce qui
sera fait au \S~9.
\end{env}

\subsection{Limites projectives de préschémas}
\label{subsection:IV.8.2-fr}

\begin{env}[8.2.1]
\label{IV.8.2.1-fr}
Soient $S_0$ un espace annelé, $L$ un ensemble préordonné filtrant croissant,
$(\mathscr A_\lambda,\varphi_{\mu\lambda})$ un système inductif de
$\mathcal O_{S_0}$-Algèbres (non nécessairement commutatives), ayant $L$ pour
ensemble d'indices. On sait que, considéré comme système inductif de
$\mathcal O_{S_0}$-Modules, $(\mathscr A_\lambda,\varphi_{\mu\lambda})$ admet
une limite inductive $\mathscr A$~; désignons par
$\varphi_\lambda:\mathscr A_\lambda\to\mathscr A$ l'homomorphisme canonique
(de $\mathcal O_{S_0}$-Modules). Soit
$m_\lambda:\mathscr A_\lambda\otimes\mathscr A_\lambda\to\mathscr A_\lambda$
l'homomorphisme de $\mathcal O_{S_0}$-Modules qui définit la multiplication
dans la $\mathcal O_{S_0}$-Algèbre $\mathscr A_\lambda$~; l'hypothèse sur les
$\varphi_{\mu\lambda}$ entraîne que les $m_\lambda$ forment un système
inductif d'homomorphismes,
